% !tp-example Basic usage
\documentclass{article}
\begin{document}
\wcDiv[data-tp-reference-scope]{
  \wcTpEquanum[ref=energy,number=2.1]{
    \rawhtml|<tp-math displaystyle="" value="E = mc^2"></tp-math>|
  }
  See equation \rawhtml*|<tp-ref href="=energy"></tp-ref>|.

}
\end{document}

% !tp-example Javascript API
\documentclass{article}
\begin{document}
\wcScript[type=module]|
import "/components/equanum/equanum.js";
import "/components/math/math.js";
import "/components/ref/ref.js";

const example = document.createDocumentFragment();

const element1 = document.createElement("div");
element1.setAttribute("data-tp-reference-scope", "");
element1.append(document.createTextNode("\n  "));

const element2 = document.createElement("tp-equanum");
element2.setAttribute("ref", "energy");
element2.setAttribute("number", "2.1");
element2.append(document.createTextNode("\n    "));

const element3 = document.createElement("tp-math");
element3.setAttribute("displaystyle", "");
element3.setAttribute("value", "E = mc^2");
element2.append(element3);
element2.append(document.createTextNode("\n  "));
element1.append(element2);
element1.append(document.createTextNode("\n  "));

const element4 = document.createElement("p");
element4.append(document.createTextNode("See equation "));

const element5 = document.createElement("tp-ref");
element5.setAttribute("href", "=energy");
element4.append(element5);
element4.append(document.createTextNode("."));
element1.append(element4);
element1.append(document.createTextNode("\n"));
example.append(element1);

document.body.append(example);
|
\end{document}

% !tp-example align row previews
\documentclass{article}
\begin{document}
\wcTpLatex{
  \wcScript[type=tp/latex]|
\usepackage{amsmath}
\begin{align}
x &= 1\label{row:first} \\
y &= 2\label{row:second}
\end{align}
Compare \eqref{row:first} and \eqref{row:second}: each reference previews only its own equation.
        |
}
\end{document}

% !tp-example alignat row previews
\documentclass{article}
\begin{document}
\wcTpLatex{
  \wcScript[type=tp/latex]|
\usepackage{amsmath}
\begin{alignat}{2}
x &= 1 &\qquad z &= 3\label{row:first} \\
y &= 2 &\qquad w &= 4\label{row:second}
\end{alignat}
Compare \eqref{row:first} and \eqref{row:second}: each reference previews only its own equation.
        |
}
\end{document}

% !tp-example flalign row previews
\documentclass{article}
\begin{document}
\wcTpLatex{
  \wcScript[type=tp/latex]|
\usepackage{amsmath}
\begin{flalign}
x &= 1 &&\label{row:first} \\
y &= 2 &&\label{row:second}
\end{flalign}
Compare \eqref{row:first} and \eqref{row:second}: each reference previews only its own equation.
        |
}
\end{document}

% !tp-example gather row previews
\documentclass{article}
\begin{document}
\wcTpLatex{
  \wcScript[type=tp/latex]|
\usepackage{amsmath}
\begin{gather}
x = 1\label{row:first} \\
y = 2\label{row:second}
\end{gather}
Compare \eqref{row:first} and \eqref{row:second}: each reference previews only its own equation.
        |
}
\end{document}

% !tp-example eqnarray row previews
\documentclass{article}
\begin{document}
\wcTpLatex{
  \wcScript[type=tp/latex]|
\usepackage{amsmath}
\begin{eqnarray}
x &=& 1\label{row:first} \\
y &=& 2\label{row:second}
\end{eqnarray}
Compare \eqref{row:first} and \eqref{row:second}: each reference previews only its own equation.
        |
}
\end{document}
